\begin{landscape}
\begin{table}[!h]
\begin{center}
\bgroup
\def\arraystretch{1.5}
    \begin{tabular}{| p{3cm} | p{5cm} | p{3cm} | p{6cm} | p{3cm} |}
    \hline
UC-ID & Use-Case-Name & Primärer Akteur & Ziel / Kurzbeschreibung & Zugeordnete User-Stories \\
\hline\hline

\textbf{UC-VAL-001} & Objektvorschlag validieren & Bürger:in & Prüft einen platzierten Objektvorschlag gegen das Regelwerk und liefert ein Valid/Invalid-Ergebnis. & US-B-001, US-B-002, US-B-004, US-B-005, US-B-006 \\
\hline
\textbf{UC-VAL-002} & Live-Feedback bei Platzierung & Bürger:in & Liefert während Drag/Move eine Anzeige durch rote Hinterlegung, damit eine gültige Position schneller gefunden wird. & US-B-003 \\
\hline
\textbf{UC-ADM-001} & Regeln definieren \& verwalten & Stadtplaner:in & Erstellt, ändert, aktiviert oder deaktiviert Regeln in einer formalen Syntax. & US-SP-001, US-SP-002, US-SP-003, US-SP-004, US-SP-005 \\
\hline
\textbf{UC-DEV-001} & Validierungs-API nutzen & AR-Entwickler:in & Ruft den Validator über ein typisiertes Interface auf und konsumiert das Ergebnis performant. & US-EAB-001, US-EAB-002 \\
\hline
\end{tabular}
    \egroup
    
\caption{Use Cases}
\label{tab:Use Cases}
\end{center}
\end{table}
\end{landscape}
\FloatBarrier
